\documentclass[aspectratio=169]{beamer}

\usetheme{Madrid}
\usecolortheme{seahorse}

\usepackage[utf8]{inputenc}
\usepackage{xcolor} % <-- FIX for \definecolor
\usepackage{amsmath, amsfonts, amssymb}
\usepackage{booktabs}
\usepackage{array}
\usepackage{tikz}

\definecolor{AccentBlue}{RGB}{0,92,184}
\setbeamercolor{structure}{fg=AccentBlue}
\setbeamercolor{title}{fg=white,bg=AccentBlue}
\setbeamercolor{frametitle}{fg=AccentBlue}

\setbeamertemplate{navigation symbols}{}
\setbeamertemplate{footline}[frame number]

%=========================================================
% Placeholder macro used in your slides (\placeholderbox{...})
%=========================================================
\newcommand{\placeholderbox}[2][0.82\linewidth]{%
  \begin{center}
    \fbox{\begin{minipage}[c][0.46\textheight][c]{#1}
      \centering \vspace{0.3em}
      \textbf{#2}\\[0.6em]
      \emph{Insert figure, diagram, or table here later.}
    \end{minipage}}
  \end{center}
}

% (Optional) Your original placeholder macro (still usable)
\newcommand{\figplaceholder}[1]{%
  \begin{center}
    \begin{tikzpicture}
      \draw[dashed] (0,0) rectangle (8,4);
      \node at (4,2) {#1};
    \end{tikzpicture}
  \end{center}
}

\title{\textbf{Inflation, Cost of Living, and Monetary Policy}}
\subtitle{ECON 204 -- Contemporary Economic Issues}
\author{Muhebullah Karimzada}
\date{Winter 2026}

\begin{document}

\begin{frame}
\transfade
\titlepage
\end{frame}

%=========================================================
\section{Inflation and the Cost of Living}

\begin{frame}
\transfade
\frametitle{What Is Inflation?}
\begin{itemize}
\item Inflation refers to a sustained increase in the general price level of goods and services over time.
\item When inflation occurs, the purchasing power of money falls, meaning households can buy fewer goods and services with the same income.
\item A one-time price increase does not necessarily constitute inflation unless it leads to continuing price growth across many sectors.
\item Inflation is therefore a macroeconomic phenomenon rather than a change in the price of a single good.
\end{itemize}
\end{frame}

\begin{frame}
\transfade
\frametitle{The Cost of Living}
\begin{itemize}
\item The cost of living refers to the amount of income required to maintain a given standard of living.
\item Economists measure changes in the cost of living using price indices.
\item The most important of these indices is the Consumer Price Index, or CPI.
\item The CPI tracks the cost of a fixed basket of goods and services typically purchased by households.
\end{itemize}
\end{frame}

\begin{frame}
\transfade
\frametitle{How the CPI Is Constructed}
\begin{itemize}
\item The CPI compares the cost of a representative basket of goods and services in the current period to its cost in a base year.
\item The basket includes items such as food, housing, transportation, clothing, and services.
\item When the CPI rises, it indicates that average consumer prices have increased.
\item This increase reflects a higher cost of living for households.
\end{itemize}
\end{frame}

\begin{frame}
\transfade
\frametitle{Measuring Inflation}
The inflation rate is calculated using the CPI formula:
\[
\pi_t = \frac{CPI_t - CPI_{t-1}}{CPI_{t-1}} \times 100
\]
\begin{itemize}
\item This formula measures the percentage change in consumer prices from one period to the next.
\item For example, if the CPI rises from 158 to 163, inflation is approximately 3.16 percent.
\item This means that average prices increased by just over three percent.
\end{itemize}
\end{frame}

\begin{frame}
\transfade
\frametitle{Why CPI Is Not a Perfect Measure}
\begin{itemize}
\item The CPI may overstate inflation because it assumes households do not substitute toward cheaper goods.
\item It may not fully account for improvements in product quality.
\item New goods and services are not always captured immediately.
\item Despite these limitations, the CPI remains the standard measure of inflation.
\end{itemize}
\end{frame}



%=========================================================
\section{Types and Causes of Inflation}

\begin{frame}
\transfade
\frametitle{Demand-Pull Inflation}
\begin{itemize}
\item Demand-pull inflation occurs when aggregate demand grows faster than productive capacity.
\item Strong household spending, business investment, and government spending can increase total demand.
\item When the economy is near full capacity, firms respond by raising prices rather than output.
\item This leads to sustained increases in the price level.
\end{itemize}
\end{frame}






\begin{frame}
\transfade
\frametitle{Cost-Push Inflation}
\begin{itemize}
\item Cost-push inflation results from rising production costs.
\item Higher wages, energy prices, or supply disruptions raise firms' costs.
\item Firms increase prices to maintain profitability.
\item Inflation occurs alongside slower economic growth.
\end{itemize}
\end{frame}





\begin{frame}
\transfade
\frametitle{Inflation Expectations}
\begin{itemize}
\item Inflation can persist when households and firms expect higher future inflation.
\item Workers demand higher wages to protect purchasing power.
\item Firms raise prices to cover higher labor costs.
\item This wage-price spiral reinforces inflation over time.
\end{itemize}
\end{frame}



%=========================================================
\section{Inflation in the AD--AS Model}

\begin{frame}
\transfade
\frametitle{Aggregate Demand and Supply}
\begin{itemize}
\item Aggregate demand represents total spending in the economy.
\item Aggregate supply represents total output firms are willing to produce.
\item The intersection of AD and AS determines the price level and real GDP.
\end{itemize}
\end{frame}

\begin{frame}
\transfade
\frametitle{AD--AS Diagrams}
\centering
\vspace{2em}

\includegraphics[width=0.7\linewidth]{Fig1}.
\end{frame}

\begin{frame}
\transfade
\frametitle{Demand-Pull Inflation in AD--AS}
\begin{itemize}
\item An increase in aggregate demand shifts the AD curve to the right.
\item In the short run, both output and prices rise.
\item Over time, rising costs shift short-run aggregate supply leftward.
\item Output returns to potential while the price level remains higher.
\end{itemize}
\end{frame}
\begin{frame}
\transfade
\frametitle{Demand-Pull Inflation}
\centering


\includegraphics[width=0.65\linewidth]{Fig3}.
\end{frame}
\begin{frame}
\transfade
\frametitle{Cost-Push Inflation in AD--AS}
\begin{itemize}
\item Rising production costs shift short-run aggregate supply leftward.
\item Prices rise while output falls.
\item Policymakers face a trade-off between inflation control and growth.
\end{itemize}
\end{frame}

\begin{frame}
\transfade
\frametitle{Cost-Push Inflation}
\centering


\includegraphics[width=0.75\linewidth]{Fig4}.
\end{frame}

%=========================================================
\section{Consequences of Inflation}

\begin{frame}
\transfade
\frametitle{Economic Costs of Inflation}
\begin{itemize}
\item High inflation creates uncertainty for households and firms.
\item Planning long-term contracts becomes more difficult.
\item Inflation redistributes income in arbitrary ways.
\item Fixed-income households may lose purchasing power.
\end{itemize}
\end{frame}

\begin{frame}
\transfade
\frametitle{Borrowers and Lenders}
\begin{itemize}
\item Unexpected inflation benefits borrowers.
\item Loan repayments are made in money with lower purchasing power.
\item Lenders receive less real value than expected.
\end{itemize}
\end{frame}

%=========================================================
\section{Nominal and Real Variables}

\begin{frame}
\transfade
\frametitle{Nominal vs Real Values}
\begin{itemize}
\item Nominal values are measured in current dollars.
\item Real values are adjusted for inflation.
\item A nominal wage increase does not guarantee higher purchasing power.
\end{itemize}
\end{frame}

\begin{frame}
\transfade
\frametitle{Real Interest Rates}
\[
r \approx i - \pi^e
\]
\begin{itemize}
\item The real interest rate reflects the true cost of borrowing.
\item It depends on expected inflation.
\item It influences consumption, saving, and investment decisions.
\end{itemize}
\end{frame}

%=========================================================
\section{Monetary Policy in Canada}

\begin{frame}
\transfade
\frametitle{The Bank of Canada}
\begin{itemize}
\item Monetary policy in Canada is conducted by the Bank of Canada.
\item The main objective is price stability.
\item Inflation targeting provides a clear policy framework.
\end{itemize}
\end{frame}

\begin{frame}
\transfade
\frametitle{Policy Interest Rate}
\begin{itemize}
\item The Bank sets a target for the overnight interest rate.
\item This influences borrowing costs throughout the economy.
\item Commercial banks adjust their lending rates accordingly.
\end{itemize}
\end{frame}

\begin{frame}
\transfade
\frametitle{Transmission Mechanism}
\begin{itemize}
\item Higher rates reduce household consumption.
\item Business investment declines.
\item Housing demand slows.
\item The exchange rate may appreciate.
\end{itemize}
\end{frame}

\begin{frame}
\transfade
\frametitle{Monetary Policy}
\centering
\vspace{2em}

\includegraphics[width=0.99\linewidth]{Fig2}.
\end{frame}

%=========================================================
\section{Inflation Expectations and Credibility}

\begin{frame}
\transfade
\frametitle{Why Expectations Matter}
\begin{itemize}
\item Stable inflation expectations support price stability.
\item Firms and workers adjust prices and wages accordingly.
\item Policy credibility reduces the need for aggressive rate changes.
\end{itemize}
\end{frame}

%=========================================================
\section{Conclusion}

\begin{frame}
\transfade
\frametitle{Key Takeaways}
\begin{itemize}
\item Inflation reduces purchasing power over time.
\item The CPI measures changes in the cost of living.
\item Monetary policy influences inflation through interest rates.
\item Expectations play a crucial role in stabilizing the economy.
\end{itemize}
\end{frame}

\end{document}
